\cotefile{90}
\scanpage{14}
\deleted{\uline{Proposition} Soit $A$ anneau de Dedekind, $G$ gpe
\inserted{plat} affine \inserted{de t.f.} sur $A$, d'algèbre $B$ satisfaisant aux conditions
du Lemme 1. \deleted{Alors} Soit $N$ comme dans le Lemme 2,
et \deleted{\unsure{la}} considérons la représentation de $G$ dans $N$ induite
par la \unsure{repr.|représ.} régulière gauche. Alors l'hom.
$G\longrightarrow\underline{\mathrm{Aut}}(N)$ est un morphisme \unsure{fidèlem$^{\mathrm t}$|fidèlement} plat}

\[
\left\{\begin{aligned}
 &\underline{B}=\varinjlim B_i\\
1)\quad &B/B_i\ \text{plat}\quad\text{\deleted{\unsure{donc|dans}}}\qquad B_i\otimes B_i\subset B\otimes B\\
2)\quad &\Delta B_i\subset B_i\otimes B_i\\
3)\quad &B_i\ \text{loc libre de rt. f.}\quad \deleted{\text{\unsure{Suppos.|Supposons} $X$ qu. cpt qu. sép.}}
\end{aligned}\right.
\]

Soit $E$ un $\underline{O}_{X}G$-module qu. coh\unsure{.}\ \deleted{\illegible{2}}

Alors 1) Pour tt sous-module qu. coh. $F$
il existe un plus petit ss-module qu.
coh. $F'$ de $F$ \deleted{tel que} stable par $G$, et
sa formation commute à tt \unsure{changem$^{\mathrm t}$|changement} de
base plat. Si $F$ est de t.f. $F'$ itou

2) Si $X$ quasi-compact quasi-séparé, alors
$E$ est $\varinjlim$ de sous-modules \inserted{qu. coh.} de t.f.
qui sont stables par $G$.

3) \inserted{\deleted{\unsure{Si} $X$ \unsure{qu. cpt qu. sép.}} Si $X$ qu. cpt qu. sép.} Soit $F\unsure{\longrightarrow|-}E$ un hom, avec $F$ quasi-coh \inserted{de t.f.}
\unsure{ou|de} d'image \unsure{engendrant|engdr|engendre} $E$. Alors il existe
% archive sheet number: 13
\scanpage{15}
\mspage{\handwritten{- 4 -}}
quant (4.5.19.1) , il suffit de montrer que $q_1^{-1}(\handwritten{\Gamma}')=q_2^{-1}(\handwritten{\Gamma}')$ . Mais
c'est une propriété \deleted{\illegible{2}} qui est vraie si \deleted{\illegible{1}} et seulement
si elle l'est après chaque changement de base $\operatorname{Spec}(\handwritten{\kappa}(x))\handwritten{\longrightarrow}X$ , où
$x$ parcourt $X$ . Autrement dit , il suffit de prouver le corollaire lorsque
$X$ est le \uline{spectre d'un corps} ; mais comme tout $X$-préschéma étale es\handwritten{t}
alors automatiquement \uline{séparé} sur $X$ (17.6.2,c')) , on est ramené au cas
envisagé au d\deleted{\unsure{e|é}}ébut de la démonstration .
% archive sheet number: 14
\scanpage{16}
un $G$-hom $\deleted{\unsure{E|F}}\,F'\unsure{-|\to}E$ surjectif, avec
$F'$ de la forme $F\otimes\check{B}_i$ -- \unsure{donc|d'où} $F'$ loc.
libre (de t.f.) si $F$ l'est.

4) Supposons $X$ qu. cpt qu. sép. et que
tt module qu. coh. de t.f. sur $X$
soit quotient d'un loc. libre de t.f.
(p. ex $X$ admet faisceau inv. ample,
ou $X$ \unsure{noeth.} régulier). Alors tt
$\unsure{\mathcal{O}|\mathcal{E}|\underline{O}}_X$-$\unsure{G|L}$-module sur $X$ est $\simeq$ \unsure{quotient} d'un
loc. libre de t.f.
% archive sheet number: 15
\scanpage{17}
\mspage{- IV-\handwritten{962} -}
que de $k$ ; il est alors immédiat que $K\otimes_k\Omega$ est composé direct d'un nombre
fini de corps , tous isomorphes à la clôture séparable $K\Omega$ de $K$ . D'autre
part , l'anneau $ {}^{\mathrm{hs}}A$ étant hensélien , $C\otimes_A({}^{\mathrm{hs}}A)$ est composé direct des annea\handwritten{ux}
locaux $(C\otimes_A({}^{\mathrm{hs}}A))_{\underline{q}_i}$ , où $\underline{q}_i$ $(1\leq i\leq s)$ parcourt les idéaux maximaux de \deleted{\illegible{1}}
$C\otimes_A({}^{\mathrm{hs}}A)$ , et ces anneaux locaux sont henséliens (18.5.10) et ont $K\Omega$ pour
corps résiduel , donc sont strictement locaux (18.8.1) . Posons pour abréger
$C'=C\otimes_A({}^{\mathrm{hs}}A)$ et $C'_i=C'_{\underline{q}_i}$ , qui est donc un anneau quotient de $C'$ ; soient
$f:C\to C'$ , $g:{}^{\mathrm{hs}}A\to C'$ , $\handwritten{\varphi}_i:C'\to C'_i$ les applications canoniques .

\diagram{\handwritten{\begin{tikzcd}[ampersand replacement=\&,column sep=small,row sep=small]
\& \& \& \handwritten{{}^{\mathrm{hs}}C}\arrow[dl,bend right=10,"\handwritten{w'}"']\\
\handwritten{C}\arrow[r,"\handwritten{f}"']\arrow[urrr,"\handwritten{v}"] \& \handwritten{C'}\arrow[r] \& \handwritten{C'_i}\arrow[ur,bend right=10,"\handwritten{w}"'] \& \\
\handwritten{A}\arrow[u]\arrow[r,"\unsure{\prime|/}"] \& \handwritten{{}^{\mathrm{hs}}A}\arrow[u,-]\arrow[uurr,"\handwritten{u}"'] \& \&
\end{tikzcd}}}

Soit $v:$ \deleted{\illegible{1}} \inserted{\handwritten{$C\to{}^{\mathrm{hs}}C$}} l'homomorphisme canonique ; il existe \deleted{\illegible{1}} un homomorphisme
local et un seul $u:{}^{\mathrm{hs}}A\to{}^{\mathrm{hs}}C$ qui , par passage aux quotients,donne l'injection
canonique $\handwritten{\Omega}\to K\Omega$ \deleted{\illegible{1}} et tel que le composé $A\longrightarrow C\xrightarrow{v}{}^{\mathrm{hs}}C$
soit égal au composé $A\longrightarrow{}^{\mathrm{hs}}A\xrightarrow{u}{}^{\mathrm{hs}}C$ (18.8.5) ; donc il existe un unique
homomorphisme local $w_{\unsure{o|0}}:C'\longrightarrow{}^{\mathrm{hs}}C$ tel que $w_o\handwritten{\circ g=u}$, $w_o\handwritten{\circ f=v}$, et puisque $ {}^{\mathrm{hs}}C$
est un anneau local , cet homomorphisme se factorise en $C'\xrightarrow{\handwritten{\varphi_i}}C'_i\xrightarrow{w}{}^{\mathrm{hs}}C$
pour un indice $i$ bien déterminé, $w$ étant un homomorphisme local . D'autre
part , l'anneau $C'_i$ étant strictement local , il existe un homomorphisme local
$w':{}^{\mathrm{hs}}C\to C'_i$ qui , par passage aux quotients , donne l'automorphisme
identique de $K\Omega$ et qui est tel que $w'\circ v=\handwritten{\varphi}_i\circ f$ \inserted{(18.8.5)} .\deleted{\illegible{1}} On déduit d'abord de là
que $w\circ w'$ est un endomorphisme de $ {}^{\mathrm{hs}}C$ qui , par passage aux quotients \deleted{\unsure{à|a}},donne
l'automorphisme identique de $K\Omega$, donc est l'identité (18.8.5) . D'autre
part , , on a $w'\circ w\circ\handwritten{\varphi}_i\circ f=w'\circ v=\handwritten{\varphi}_i\circ f$ , et $w'\circ w\circ\handwritten{\varphi}_i\circ g=w'\circ u=\handwritten{\varphi}_i\circ g$ ; d'où $w'\circ w\circ\handwritten{\varphi}_i=
\handwritten{\varphi}_i$ , autrement dit $w'\circ w$ est l'automorphisme identique de $C'_i$ . CQFD .

\uline{Remarque} (18.8.{\deleted{7}\handwritten{10}}.1) .- La démonstration de (18.8.7)
% slanted
\handwritten{\unsure{!}}\inserted{\handwritten{! pour une $A$-algèbre $B$ \uline{entière} sur $A$}} n'ut\uline{i}lise le
fait que $B$ est une $A$-algèbre finie,que pour établir que $K\otimes_k\Omega$ est
composé direct d'un nombre fini de corps . Or , cette dernière propr\handwritten{i}été
est encore vérifiée lorsque l'on suppose que $B$ est une $A$-algèbre
\uline{entière semi-locale} dont les idéaux maximaux $\underline{n}_j$ ont des corps résiduels
de degrés séparables $\handwritten{[\kappa}(\underline{n}_j):k\handwritten{]}_s$ \uline{finis} .

\uline{Proposition} (18.8.\deleted{8}\handwritten{11}) .- \uline{Soit} $A$ \uline{un anneau} \deleted{\illegible{1}} \uline{semi-lo\unsure{d|c}al} .

(i) \uline{Pour que} $ {}^{\mathrm{hs}}A$ \uline{soit réduit} (resp. \uline{normal) , il faut et il suffit que} $A$ \uline{le soit} .

(ii) \uline{Supposons} $A$ \uline{noethérien . Alors , pour tout idéal premier} $\underline{p}$ \uline{de} $A$ ,\uline{l'anneau} $({}^{\mathrm{hs}}A)_{\underline{p}}/\underline{p}.({}^{\mathrm{hs}}A)_{\underline{p}}$ \uline{est composé direct d'un nombre fini d'extensions \unsure{algé|algè}-}
% archive sheet number: 16
\scanpage{18}
$G$ un groupe plat et affine sur $S$,
d'algèbre $B$. On suppose
\[
B=\varinjlim B_i,\qquad B_i\subset B\quad\text{sous-Modules avec}
\]
\[
\begin{aligned}
1)\quad &B/B_i\ \text{plat}\qquad(\text{donc }B_i\otimes B_i\subset B\otimes B)\\
2)\quad &\Delta B_i\subset B_i\otimes B_i\qquad\text{donc }\check B_i\text{ est un faisceau}\\
&\hspace{12em}\text{d'algèbres associatives à unité}\\
3)\quad &B_i\quad\text{\deleted{projectif}\inserted{loc libre} de t.f.}\qquad\text{donc les }\check B_i\text{ itou \unsure{!|?}.}
\end{aligned}
\]

$X$ un préschéma sur $S$, \deleted{\unsure{muni} de $\underline L$ \uline{ample} rel/$S$} \inserted{où $G$ opère.}

\deleted{$G$ \unsure{opè|op.|roupe}} On suppose que sur $X$ existe un
faisceau \inserted{\unsure{inv.}} $\underline L$ ample rel/$S$, sur lequel $G$
opère de façon compatible avec son action sur $X$.

Soit $E$ un $\unsure{\mathcal{O}|\mathcal{E}|\underline{O}}_X$-$G$-Module qu. coh. \deleted{\unsure{sur}
\unsure{lequel $G$ opè}}

1) Pour tt sous-module qu coh $F$ de $E$,
il existe un plus petit sous-module qu.
coh $F'$ de $E$ \deleted{qui} stable par $G$ \unsure{cont$^{\mathrm{nt}}$|contenant} $F$.
Sa \deleted{\unsure{changemt} de} formation commute à
tt changement de base plat en $S$.
\inserted{Si $F$ est de t.f., de \unsure{\=m|même} $F'$.}

\margin{\deleted{\unsure{t.f.|de t.f.}}}
\deleted{2) Si $S$ qu. cpt qu. sép., alors $E$ est
$\varinjlim$ de sous-modules qu. coh. de t.f.
qui sont \uline{stables} par $G$.}
% archive sheet number: 17
\scanpage{19}
\mspage{- IV-\handwritten{953} -}
En effet , $Y$ étant compact \deleted{\illegible{1}} et noethérier , toute partie de $Y$ est compacte\handwritten{,}
donc fermée , ce qui implique que $Y$ est discret , donc fini puisqu'il est
compact .

\uline{Co\deleted{\unsure{b|o}}rollaire} (18.6.10) .- \uline{Soit} $A$ \uline{un anneau local noethérien . Pour que} $ {}^{\mathrm h}A$ \uline{poss\handwritten{è}de une des propriétés suivantes} :

a) \uline{être un anneau de Cohen-Macaulay} ;

b) \uline{vérifier la propriété} $(S_n)$ ;
\margin{\handwritten{\unsure{virg.|ví}}}

c) \uline{être régulier} ;

d) \uline{vérifier la propriété} $(R_n)$ ;

\uline{il faut et il suffit que} $A$ \uline{possède cette même propriété} .

Cela résulte aussitôt de ce que les fibres du morphisme $\operatorname{Spec}({}^{\mathrm h}A)\to\operatorname{Spec}(A)$
sont géométriquement régulières (18.6.\unsure{9|8}) et de (6.4.1),(6.5.1) et (6.5.3).

\uline{Corollaire} (18.6.11) .- \uline{Soit} $A$ \uline{un anneau semi-local noethérien . Pour qu'un
idéal premier} $\underline p$ \uline{de} $ {}^{\mathrm h}A$ \uline{appartienne à} $\operatorname{Ass}({}^{\mathrm h}A)$ , \uline{il faut et il suffit que} $\underline p\cap A$
\uline{appartienne à} $\operatorname{Ass}(A)$ .

Compte tenu de la description des fibres du morphisme $\operatorname{Spec}({}^{\mathrm h}A)\to\operatorname{Spec}(A)$ ,
cela résulte aussitôt de (3.3.1) .

\uline{Proposition} (18.6.\unsure{12}) .- \uline{Soi\deleted{en}t} $A$ \uline{un anneau local ; les propriétés suivantes
sont équivalentes} :

a) $A$ \uline{est unibranche}\margin{\handwritten{\unsure{,,|virg.}}} (resp. \uline{réduit et unibranche}) .

b) \uline{Pour toute} $A$-\uline{algèbre locale strictement essentiellement étale} $A'$ ,\deleted{\illegible{1}}
$\operatorname{Spec}(A')$ \uline{est irréductible} (resp. \uline{intègre}) .

c) $\operatorname{Spec}({}^{\mathrm h}A)$ \uline{est irréductible} (resp.\uline{intègre}) .

Notons d'abord que $ {}^{\mathrm h}(A_{\mathrm{red}})=({}^{\mathrm h}A)_{\mathrm{red}}$ , car \deleted{\illegible{4}
\unsure{$A_{\mathrm{red}}$} \illegible{4}} $ {}^{\mathrm h}(A_{\mathrm{red}})=({}^{\mathrm h}A)\otimes_A A_{\mathrm{red}}$ (18.6.8) , et
comme $ {}^{\mathrm h}(A_{\mathrm{red}})$ est réduit (18.6.9) , il est égal à $({}^{\mathrm h}A)_{\mathrm{red}}$ . Ceci permet ,
dans la démonstration , de considérer seulement le cas où $A$ est \uline{réduit} ,donc
aussi $A'$ (17.5.7) et $ {}^{\mathrm h}A$ (18.6.9) .

\deleted{\illegible{2}} Avec les notations de (18.6.4) , la condition b) signifie
que tous les $A_{\handwritten{\alpha}}$ sont intègres ; \deleted{\unsure{comme les homomorphismes} \unsure{$A_{\alpha}$ à $A_{\beta}$} sont
\unsure{injectifs} ,} on en conclut que $ {}^{\mathrm h}A=\varinjlim A_{\handwritten{\alpha}}$ est intègre (5.13.3 ) , donc b) entraîne
c) .
% archive sheet number: 18
\scanpage{20}
\deleted{3) Si $\underline E$ est \deleted{\unsure{de}} de t.f. \inserted{et est quotient d'un module loc libre (\unsure{sans op.!})}, il existe un
$\unsure{\mathcal O|\mathcal E|\underline{O}}_X$-$G$-Module \deleted{\unsure{qu.}} \uline{loc. libre} \uline{de t.f.} $\underline F$,
et un épimorphisme de $\mathcal O_X$-$G$-Modules
$\underline F\longrightarrow\underline E$.}

3) \deleted{Si $S$ est \unsure{p.} admet un faisceau inv.}
\inserted{qu. cpt \deleted{\unsure{qu. sép.}} qu. sép. et si tt module quasi coh.}
de t.f. sur $S$ est $\simeq$ quotient d'un loc. libre de t.f.,
\deleted{ample,} alors $\underline E$ est $\unsure{\simeq|\cong}$ quotient d'un
$\underline O_X$-$G$-Module loc. libre de type fini
de la forme $f^*(H)(-n)$, où

\uline{Dém} Il suffit de prouver 1) et 3). Posons
(\unsure{où|ici} $f:X\unsure{-|\to}S$)
\[
\underline{\Gamma}_{\deleted{\unsure{*}}}^{\unsure{\bullet|\cdot}}(E)=\coprod_n f_*(E(n))
\]
Alors grâce au fait que $G$ est plat sur $S$
et $f$ qu. cpt et (qu) séparé, on trouve que
$\Gamma^\bullet(E)$ est un Module gradué sur $S$ sur
lequel $G$ opère. Soit
\[
\underline{\Gamma}_{\deleted{*}}^\bullet(\mathcal O_X)=\deleted{\unsure{A_*|A^\bullet}}\ A^\bullet
\]
c'est une Algèbre graduée sur \unsure{lesquels|laquelle}
$G$ opère. $E\unsure{\leadsto|\sim|\mapsto}\underline{\Gamma}^{\cdot}(E)$ est un foncteur en
$E$, à valeurs dans les \deleted{\unsure{objets}}
$A^{\cdot}$-$G$-Modules gradués. Ce foncteur
est exact à gauche. On a un foncteur
(\uline{exact})
% archive sheet number: 19
\scanpage{21}
\mspage{- IV-\handwritten{\unsure{952|953}} -}
\uline{neau} $({}^{\mathrm h}A)_{\underline p}/\underline p.({}^{\mathrm h}A)_{\underline p}$ \uline{est} \deleted{\illegible{1}} \uline{composé direct d'un nombre fini d'extensions
algébriques séparables de} $\handwritten{\kappa}(\underline p)$ (ce qui entraîne que les fibres du morphisme
canonique $\operatorname{Spec}({}^{\mathrm h}A)\to\operatorname{Spec}(A)$ sont \uline{géométriquement régulières}) .

(i) Comme $ {}^{\mathrm h}A$ est un $A$-module fidèlement plat,il résulte de (2.1.\unsure{13}) que si
$ {}^{\mathrm h}A$ est \deleted{\illegible{1}} réduit (resp. normal) , il en est de même de $A$ . Pour prouver la
réciproque , on peut se borner au cas où l'anneau $A$ est local , de sorte \deleted{\illegible{1}}
qu'avec les notations de (18.6.4) , on a $ {}^{\mathrm h}A=\varinjlim A_{\handwritten{\alpha}}$ . Or , il résulte de la
définition des $A_{\handwritten{\alpha}}$ et de (17.5.7) que si $A$ est réduit (resp. intègre et intégralement
clos) , il en est de même des $A_{\handwritten{\alpha}}$ ; \deleted{\unsure{xx}} en outre , les \deleted{\illegible{1}
\illegible{5}} morphismes \deleted{$\operatorname{Spec}(A_{\handwritten{\beta}})\to\operatorname{Spec}(A_{\handwritten{\alpha}})$} \inserted{\handwritten{$A_\alpha\unsure{-|\to}A_\beta$}} pour $\handwritten{\alpha\leq\beta}$
% slanted
\margin{\handwritten{\unsure{pinceaux|précisions} \unsure{mélangés|mélangées}!}}
sont \deleted{surjectifs} \inserted{\handwritten{injectifs}} en vertu de \deleted{(2.5.4,(ii))}\inserted{\handwritten{\unsure{----}}} . On conclut donc de (5.13.2 ) (resp
(5.13.4 )) que $ {}^{\mathrm h}A$ est réduit (resp. intègre et intégralement clos) .

(ii) \inserted{On peut se borner au cas où $A$ est local .} En vertu du fait que le foncteur $\varinjlim$ permute au produit tensoriel ,
la fibre du morphisme $\operatorname{Spec}({}^{\mathrm h}A)\to\operatorname{Spec}(A)$ en un point $p$ est limite inductive
des fibres des morphismes $\operatorname{Spec}(A_{\handwritten{\alpha}})\to\operatorname{Spec}(A)$ en ce point . . Comme $ {}^{\mathrm h}A$ est
noethérien , on voit , compte tenu de (17.6.2) , que l'on est ramené à prouver
le lemme suivant :

\uline{Lemme} (18.6.9.1) .- \uline{Soient} \deleted{\illegible{1}} $(B_{\handwritten{\alpha}})$ \uline{un système inductif d'anneaux artiniens} ,
$B=\varinjlim B_{\handwritten{\alpha}}$ \uline{sa limite inductive . Si} $B$ \uline{est noethérien , alors} $B$ \uline{est artinien} ;
\deleted{\illegible{6}} \uline{si de plus , pour} $\handwritten{\alpha\leq\beta}$ , \uline{les
homomorphismes} \inserted{\handwritten{$\varphi_{\beta\alpha}:$}} $B_{\handwritten{\alpha}}\to B_{\handwritten{\beta}}$ \uline{sont injectifs , alors il existe} $\handwritten{\lambda}$ \uline{tel que , pour} $\handwritten{\alpha\geq\lambda}$
\uline{le nombre des composants locaux} \deleted{\illegible{1}} $B_{\handwritten{\alpha}}^{(i)}$ \uline{de} $B_{\handwritten{\alpha}}$ \uline{soit constant} , \deleted{\illegible{1}}
\uline{que} $\handwritten{\varphi_{\beta\alpha}}(B_{\handwritten{\alpha}}^{(i)})\subset B_{\handwritten{\beta}}^{(i)}$ \uline{pour tout $i$ et que les} $B^{(i)}=\varinjlim B_{\handwritten{\alpha}}^{(i)}$ \uline{soient les composant\handwritten{s}
locaux de} $B$ .

Soient en effet $Y_{\handwritten{\alpha}}=\operatorname{Spec}(B_{\handwritten{\alpha}})$ , $Y=\operatorname{Spec}(B)$ , qui , en tant qu'espace topologique
, est égal à $\varprojlim Y_{\handwritten{\alpha}}$ (8.2.10) . Les espaces $Y$ étant \deleted{\illegible{1}} \inserted{\handwritten{finis et}} discrets , la seconde
assertion du lemme résultera de la première et du fait que si les $\handwritten{\varphi_{\beta\alpha}}$
% slanted
\margin{\handwritten{\unsure{pas très clair}}}
sont injectifs , les morphismes canoniques $Y_{\handwritten{\beta}}\to Y_{\handwritten{\alpha}}$ sont surjectifs . On est
ainsi ramené à prouver le lemme purement topologique suivant :

\uline{Lemme} (18.6.9.2) .-\uline{Soit} $(Y_{\handwritten{\alpha}})$ \uline{un système projectif d'espaces topologiques finis
et discrets . Si} $Y=\varprojlim Y_{\handwritten{\alpha}}$ \uline{est noethérien} , $Y$ \uline{est fini et discret} .
% archive sheet number: 20
\scanpage{22}
$M_{\cdot}\unsure{\leadsto|\sim|\mapsto}\widetilde M_{\cdot}$ en sens inverse, \inserted{(ils sont adjoints \unsure{---|...})} et le composé \deleted{est}
$E\leadsto\widetilde{\underline{\Gamma}^{\unsure{\cdot|\bullet}}(E)}$ est isom. à l'identité, ce qui
montre que $\underline{\Gamma}^{\cdot}$ est \uline{pleinement fidèle}.

$F$ correspond à un sous-\deleted{\unsure{lin|G}}$A^{\cdot}$-module
$\Gamma^{\cdot}(F)$ de $\Gamma^{\cdot}(E)$, pas à priori stable par action
de $G$. Si \deleted{$F'\supset F$ est} \inserted{$F'$ tel que} $F\subset F'\subset E$ \deleted{\unsure{et}}
\deleted{$F'\supset F$,} est stable par $G$, alors $\Gamma^{\cdot}(F')$
\deleted{$\Gamma_{\deleted{\unsure{*}}}^{\cdot}(E')\unsure{\subset|\supset}\Gamma^{\cdot}(E)$,} satisfait
$\underline{\Gamma}^{\cdot}(F)\subset\underline{\Gamma}^{\cdot}\unsure{[|(}F')\subset\underline{\Gamma}^{\cdot}(E)$ et est stable par $G$.
Or on sait qu'il y a un plus petit
sous-\inserted{$\mathcal O_S$-}Module $N$ de $\Gamma^{\cdot}(E)$, \unsure{cont$^{\mathrm{nt}}$|contenant}
$\underline{\Gamma}^{\cdot}(F)$, et stable par $G$. On voit tout de
suite que $N$ est un sous-module gradué
[ \unsure{par|car|pour} \unsure{chaque} $N\cap\Gamma^i(E)$ \unsure{contient} $\underline{\Gamma}^i(F)$, et
est stable par $G$, donc contient le plus
petit ss-module $M^i$ de $\underline{\Gamma}^i(E)$ \unsure{contenant}
$\underline{\Gamma}^i(F)$ et stable par $G$, donc $N$ contient
$\sum M^i$, donc par minimalité lui est
égal \unsure{]|$\}$|?]}. De \unsure{\=m|même}, par fonctorialité de la
construction de $N$, on voit que $N$ est
un sous-$A^\cdot$-module. Considérons alors
% archive sheet number: 21
\scanpage{23}
\mspage{- IV-\handwritten{951} -}
et il est clair p\deleted{\unsure{q}}\unsure{a|q}r (18.\unsure{6}.6,(v)) que si $A$ est n\unsure{o|h}ethérien \deleted{\illegible{1}} , il en est de
même de $ {}^{\mathrm h}A$ . \deleted{Enfin , $ {}^{\mathrm h}A$ possède aussi la propriété universelle (18.6.6,(ii))
en convenant d'appeler``homomorphisme \inserted{\handwritten{semi-}}local'' de $A$ dans un anneau \inserted{\handwritten{semi-}}\uline{local} $B$ tout
homomorphisme $\handwritten{\varphi}$ tel que $\handwritten{\varphi}(\underline r)$ soit contenu dans l'idéal maximal \inserted{$\underline n$} de $B$ . En
effet , il est clair que pour un tel homomorphisme , $\handwritten{\varphi}^{-1}(\underline n)$ doit contenir
l'un des $\underline m_i$ puisqu'il contient leur intersection , donc $\handwritten{\varphi}^{-1}(\underline n)$ est égal à
l'un des $\underline m_i$ , et $\handwritten{\varphi}$ se factorise \deleted{\unsure{donc}} en $A\to A_{\underline m_i}\to B$ ; d'où l'on conclut
\deleted{\illegible{1}} que $\operatorname{Hom.loc}(A,B)=\handwritten{\prod_i}\operatorname{Hom.loc}(A_{\underline m_i},B)$ , et le fait que (18.6.6,(ii))
est encore valable lorsque $A$ est semi-local en découle aussitôt .}
% slanted
\margin{\handwritten{vider, (la forme hybride ici ne mérite pas le nom de ``\unsure{prop|propr.} universelle'') ou bien \unsure{arranger} avec \deleted{des} ``hom $1/2$ locaux'' dans \deleted{des} anneaux \unsure{henséliens}}}

\uline{Proposition} (18.6.8) .- \uline{Soient} $A$ \uline{un anneau semi-local} , $B$ \uline{une} $A$-\uline{algèbre semi-locale et entière sur} $A$ . \uline{Alors} $B\otimes_A({}^{\mathrm h}A)$ \uline{est une} $B$-\uline{algèbre semi-locale isomorphe à} $ {}^{\mathrm h}B$ .

\deleted{\illegible{7}} Comme $B\otimes_A({}^{\mathrm h}A)$ est
entier au-dessus de $ {}^{\mathrm h}A$ , chacun de ses idéaux maximaux est au-dessus d'un
des idéaux maximaux de $ {}^{\mathrm h}A$ , donc au-dessus d'un idéal maximal de $A$ , et comme
$B$ est entière sur $A$ , les idéaux premiers de $B$ au-dessus d'un idéal maximal
de $A$ sont des idéaux maximaux de $B$ , si bien que finalement un idéal maximal
de $B\otimes_A({}^{\mathrm h}A)$ a pour projections dans $\operatorname{Spec}({}^{\mathrm h}A)$ et $\operatorname{Spec}(B)$ des idéaux maximaux
; compte tenu de (18.6.6,(iii)) et de (\uline{I},3.4.9) , on conclut que l'anneau
$B\otimes_A({}^{\mathrm h}A)$ est semi-local . En outre , pour tout anneau local hensélien\inserted{$C$},
$\operatorname{Hom.loc}(B\otimes_A({}^{\mathrm h}A),C)$ est en correspondance biunivoque avec l'ensemble des
couples $(\handwritten{\varphi},\handwritten{\psi})$ d'homomorphismes locaux $\handwritten{\varphi}:B\to C$ , $\handwritten{\psi}:{}^{\mathrm h}A\to C$ tels que les composés
$A\to B\xrightarrow{\handwritten{\varphi}}C$ , $A\to{}^{\mathrm h}A\xrightarrow{\handwritten{\psi}}C$ soient égaux . Mais \deleted{\illegible{2}} en vertu
de la correspondance biunivoque entre $\operatorname{Hom.loc}(A,C)$ et $\operatorname{Hom.loc}({}^{\mathrm h}A,C)$ , on
voit que pour tout $\handwritten{\varphi\in}\operatorname{Hom.loc}(B,C)$ il existe un $\handwritten{\psi}$ et un seul ayant \deleted{\unsure{les}} \inserted{la} propriété\deleted{s}
précédente\deleted{s} , donc l'application $\operatorname{Hom}(B\otimes_A({}^{\mathrm h}A),C)\to\operatorname{Hom}(B,C)$ est bijective
, ce qui prouve la proposition en vertu de l'unicité de la solution
d'un problème universel .
% slanted
\margin{\handwritten{\uline{\unsure{cavalier|cavalé}}}}

\uline{Théorème} (18.6.9) .- \uline{Soit} $A$ \uline{un anneau semi-local} .

(i) \uline{Pour que} $ {}^{\mathrm h}A$ \uline{soit réduit} (resp. \uline{normal),il faut et il suffit que} $A$ \uline{le soi}\handwritten{t}.

(ii) \uline{Supposons} $A$ \uline{noethérien . Alors , pour tout idéal premier} $\underline p$ \uline{de} $A$ , \uline{l'an}-
% archive sheet number: 22
\scanpage{24}
$\widetilde N$, c'est un $\underline O_X$-$G$-module \deleted{\unsure{sur $X$} \unsure{conte|contenu}}
$\subset\inserted{\unsure{\widetilde{\underline{\Gamma}^{\cdot}(F')}|\underline{\Gamma}^{\cdot}(F')}=}\deleted{\unsure{E}}F'$, et contenant $\widetilde{\underline{\Gamma}^{\cdot}(F)}=F$. C'est donc le
\deleted{\unsure{$F'$} cherché} faisceau cherché. Par \unsure{construct$^{\mathrm n}$|construction},
sa formation commute à \unsure{toute} extension
plate $S'\to S$ de la base. Enfin, si
$F$ est de type fini, \deleted{\unsure{alors $\exists$}} prouvons que
son enveloppe $F'$ l'est. On peut supposer $S$
affine. Alors $\exists\,n$ tel que $F(n)$
\deleted{\unsure{pour} \unsure{tt} $\unsure{i\leq N}$}
engendré par ses sections, \deleted{\unsure{soient des}} \inserted{donc} par des
sections $\varphi_i\deleted{\unsure{*}}\ (1\leq i\leq N)\in\Gamma(\unsure{\mathcal O_S|S},\underline{\Gamma}^n(F))$.

Considérons alors le sous-module
\deleted{$\underline N^n$ de $\underline{\Gamma}^{\unsure{\cdot|\bullet}}(E)$ engendré par les $\varphi_i$.
Alors $\widetilde{\underline N}^{\cdot}\subset F$ est en \deleted{\unsure{un}} fait égal à $F$,
comme on constate tt de suite. On
peut alors répéter la construction précédente,}
\unsure{ou|au} $G$-invariant de $\underline{\Gamma}^{\unsure{n|-}}(F)$ engendré
par les $\varphi_i$, soit $N^n$, et le sous
Module de $\underline E$ qu'il \unsure{définit|défint}. \uline{On
vérifie} \unsure{quand même} que c'est $\underline F'$, ce qui
prouve que $F'$ est de type fini\unsure{.....|......}
% slanted
\margin{Rédiger mieux \unsure{en|au|ce} début de la démonstration, pour pouvoir y référer ici \unsure{.....|...}}
% archive sheet number: 23
\scanpage{25}
\mspage{- IV-\handwritten{950} -}
$f_{\unsure{o|0}}^{(\handwritten{\lambda})}:S_o^{(\handwritten{\lambda})}\to {S'}_o^{(\handwritten{\lambda})}$ tels que $S'=S'^{(\handwritten{\lambda})}\times_{S^{(\handwritten{\lambda})}}S$ , $g=g^{(\handwritten{\lambda})}\times 1$ et $f_o=f_o^{(\handwritten{\lambda})}\times 1$ . \deleted{\illegible{1}}
Soient $s$ le point fermé de $S$ , $x=f_o(s)\deleted{\unsure{{}_s}}$ , $x_{\handwritten{\lambda}}$ la projection de $x$ dans $S'^{(\handwritten{\lambda})}$ ;
comme $x_{\handwritten{\lambda}}$ est au-dessus du point fermé de $S^{(\handwritten{\lambda})}$ , l'anneau local \inserted{\handwritten{$C_\lambda$}} de $S'^{(\handwritten{\lambda})}$ au
point $x_{\handwritten{\lambda}}$ est une $A_{\handwritten{\lambda}}$-algèbre essentiellement étale ; en outre , comme $f_o^{(\handwritten{\lambda})}(s_{\handwritten{\lambda}})=$
$=x_{\handwritten{\lambda}}$ , on a $\handwritten{\kappa}(x_{\handwritten{\lambda}})=\handwritten{\kappa}(s_{\handwritten{\lambda}})$ , autrement dit , $\unsure{C|S|\mathcal{O}}_{\handwritten{\lambda}}$ est une $A_{\handwritten{\lambda}}$-algèbre locale strictement
essentiellement étale , et par suite (18.6.1) est $A_{\handwritten{\lambda}}$-isomorphe à une
$A$-algèbre $A_{\handwritten{\mu}}$ avec $\handwritten{\mu\geq\lambda}$ . Il y a \deleted{\illegible{2}} \inserted{\handwritten{donc}} un $A_{\handwritten{\lambda}}$-homomorphisme \deleted{\illegible{1}} \inserted{\handwritten{$C_\lambda\to{}^{\mathrm h}A$}} , c'est-à-dire
un $S^{(\handwritten{\lambda})}$-morphisme $h:S\to S'^{(\handwritten{\lambda})}$ tel que $h(s)=x_{\handwritten{\lambda}}$ , et par suite il existe
bien une $S$-section $f$ de $S'$ telle que $f(s)=x$ , ce qui achève de prouver (i) \unsure{;|.}

Pour prouver (ii) , il suffit de noter que l'on a $\operatorname{Hom.loc}({}^{\mathrm h}A,B)\deleted{\unsure{=}}$
\margin{\handwritten{$\simeq$}}
$=\varprojlim\operatorname{Hom.loc}(A_{\handwritten{\lambda}},B)$ et que par (18.6.2) les homomorphismes canoniques
$\operatorname{Hom.loc}(A_{\handwritten{\lambda}},B)\longleftarrow\operatorname{Hom.loc}(A,B)$ sont bijectifs .

\deleted{\illegible{1}} Pour démontrer \deleted{\illegible{1}} (iv) , \deleted{\illegible{2}} notons que pour tout $\handwritten{\lambda}$ et tout
entier $n>0$ , $\underline m_{\handwritten{\lambda}}^n=\underline m^n A_{\handwritten{\lambda}}$ , et $(\underline m.{}^{\mathrm h}A)^n=\underline m^n.{}^{\mathrm h}A=\varinjlim\underline m_{\handwritten{\lambda}}^n$ . On en conclut , par l'exactitude
du foncteur $\varinjlim$ dans la catégorie des $A$-modules , que $ {}^{\mathrm h}A/(\underline m.{}^{\mathrm h}A)^n=$
$=\varinjlim(A_{\handwritten{\lambda}}/\underline m_{\handwritten{\lambda}}^n)$ , et par suite , il suffit de montrer que , pour tout entier $n$,
et tout indice $\handwritten{\lambda}$ , l'homomorphisme $A/\underline m^n\to A_{\handwritten{\lambda}}/\underline m_{\handwritten{\lambda}}^n$ est bijectif . Or , \deleted{\illegible{1}} cela
est vrai par hypothèse pour $n=1$ ; d'autre part , puisque $A_{\handwritten{\lambda}}$ est un $A$-module
plat , on a
\[
\underline m_{\handwritten{\lambda}}^n/\underline m_{\handwritten{\lambda}}^{n+1}
=(\underline m^n/\underline m^{n+1})\otimes_A A_{\handwritten{\lambda}}
=(\underline m^n/\underline m^{n+1})\otimes_{A/\underline m}(A_{\handwritten{\lambda}}/\underline m A_{\handwritten{\lambda}})
\]
et comme $A/\underline m\to A_{\handwritten{\lambda}}/\underline m A_{\handwritten{\lambda}}=A_{\handwritten{\lambda}}/\underline m_{\handwritten{\lambda}}$ est bijectif , l'homomorphisme $\underline m^n/\underline m^{n+1}\to\underline m_{\handwritten{\lambda}}^n/\underline m_{\handwritten{\lambda}}^{n+1}$
est lui aussi bijectif ; la conclusion résulte donc de Bourbaki,\uline{Alg.comm.},
chap.III,§ 2,n°8,\handwritten{cor. 3 du th. 1} .

Enfin , si $A$ est hensélien \deleted{[ et si $A'=B_{\underline n}$ est une $A$-algèbre locale
strictement essentiellement étale , où $B$ est une $A$-algèbre étale , il
résulte de \deleted{\illegible{1}} (18.6.2) qu'il existe une \deleted{\illegible{1}} $\operatorname{Spec}(A)$-\uline{section de}
$\operatorname{Spec}(B)$ prenant la valeur $\underline n$ au point fermé de $\operatorname{Spec}(A)$ ]} ; on conclut
\deleted{alors} \inserted{\handwritten{de la remarque qui \unsure{s|suit}}} de \deleted{(17.4.1)} \inserted{\handwritten{18.6.2}} que \deleted{\illegible{2}} \deleted{l'}\inserted{\handwritten{les}} homomorphisme\handwritten{s} \deleted{\illegible{1}} $A\to A\deleted{'}_{\handwritten{\lambda}}$ \deleted{est} \inserted{\handwritten{sont}}
bijectif\handwritten{s}, ce qui prouve (vi) par définition de $ {}^{\mathrm h}A$ .

(18.6.7) Soient maintenant $A$ un anneau \uline{semi-local} , $\underline m_i$ $(1\leq i\leq r)$ ses idéaux
maximaux . On appelle \uline{hensélisé} de $A$ et l'on note encore $ {}^{\mathrm h}A$ l'anneau \uline{produit}
$\handwritten{\prod_i}{}^{\mathrm h}(A_{\underline m_i})$ des hensélisés des anneaux locaux $A_{\underline m_i}$ . C'est \inserted{\handwritten{un $A$-module fidèlement plat et}} une $A$-algèbre semi-locale
, dont les idéaux maximaux sont les $\underline m_i.{}^{\mathrm h}A$ en vertu de (18.\unsure{6|5}.6,(iii));
en outre , si $\underline r=\handwritten{\bigcap_i}\underline m_i$ est le radical de $A$ , il résulte de ce qui précède
que $\underline r.{}^{\mathrm h}A$ est le radical de $ {}^{\mathrm h}A$ et que l'application $A/\underline r\to{}^{\mathrm h}A/\underline r.{}^{\mathrm h}A$ est bijective
. Comme le séparé complété $\hat A$ de $A$ pour la topologie $\underline r$-préadique est le
produit des séparés complétés $\hat A_{\underline m_i}$ (Bourbaki,\uline{Alg.comm.},chap.III,§ 2,n°13 \handwritten{prop.
18}) l'homomorphisme canonique $\hat A\to({}^{\mathrm h}A)\handwritten{{}^{\wedge}}$ est bijectif par (18.6.6,(iv)),
% archive sheet number: 24
